
\documentclass{article}
\usepackage{colortbl}
\usepackage{makecell}
\usepackage{multirow}
\usepackage{supertabular}

\begin{document}

\newcounter{utterance}

\centering \large Interaction Transcript for game `matchit\_ascii', experiment `similar\_grid\_1', episode 3 with gemma{-}4{-}26b{-}a4b{-}it.
\vspace{24pt}

{ \footnotesize  \setcounter{utterance}{1}
\setlength{\tabcolsep}{0pt}
\begin{supertabular}{c@{$\;$}|p{.15\linewidth}@{}p{.15\linewidth}p{.15\linewidth}p{.15\linewidth}p{.15\linewidth}p{.15\linewidth}}
    \# & $\;$A & \multicolumn{4}{c}{Game Master} & $\;\:$B\\
    \hline

    \theutterance \stepcounter{utterance}  
    & & \multicolumn{4}{p{0.6\linewidth}}{
        \cellcolor[rgb]{0.9,0.9,0.9}{
            \makecell[{{p{\linewidth}}}]{
                \texttt{\tiny{[P1$\langle$GM]}}
                \texttt{Συμμετέχετε σε ένα συνεργατικό παιχνίδι μαντεψιάς.} \\
\texttt{Ο στόχος είναι να βρείτε αν αυτό το πλέγμα και μια άλλη εικόνα που μόνο εγώ μπορώ να δω, είναι τα ίδια.} \\
\texttt{Παρακαλώ περιγράψτε πρώτα το πλέγμα σας. Στη συνέχεια, θα δώσω τη δική μου περιγραφή και μπορούμε να κάνουμε ο ένας στον άλλο ερωτήσεις σχετικά με τα πλέγματα για να καταλάβουμε αν είναι τα ίδια. Ξεκινήστε τώρα τη σύντομη περιγραφή του πλέγματός σας με "ΠΕΡΙΓΡΑΦΗ:" ακολουθούμενη από την περιγραφή. Μην προσθέσετε τίποτα άλλο. Αυτό είναι το πλέγμα σας:} \\
\texttt{X X X X X} \\
\texttt{X ▢ ▢ ▢ ▢} \\
\texttt{X X X ▢ ▢} \\
\texttt{X ▢ ▢ ▢ ▢} \\
\texttt{X ▢ ▢ ▢ ▢} \\
            }
        }
    }
    & \\ \\

    \theutterance \stepcounter{utterance}  
    & \multicolumn{4}{p{0.6\linewidth}}{
        \cellcolor[rgb]{0.9,0.9,0.9}{
            \makecell[{{p{\linewidth}}}]{
                \texttt{\tiny{[P1$\rangle$GM]}}
                \texttt{ΠΕΡΙΓΡΑΦΗ: Ένα πλέγμα 5x5 όπου οι πρώτες στήλες και οι πρώτες γραμμές περιέχουν κυρίως "X". Συγκεκριμένα, η πρώτη στήλη είναι γεμάτη με "X". Η πρώτη γραμμή είναι γεμάτη με "X". Η δεύτερη γραμμή έχει "X" στην πρώτη θέση και μετά τέσσερα κενά (▢). Η τρίτη γραμμή έχει τρία "X" στην αρχή και δύο κενά (▢). Οι τέσσερτες και πέμπτες γραμμές έχουν ένα "X" στην αρχή και τέσσερα κενά (▢).} \\
            }
        }
    }
    & & \\ \\

    \theutterance \stepcounter{utterance}  
    & & & \multicolumn{2}{p{0.3\linewidth}}{
        \cellcolor[rgb]{0.9,0.9,0.9}{
            \makecell[{{p{\linewidth}}}]{
                \texttt{\tiny{[GM$|$GM]}}
                \texttt{continue} \\
            }
        }
    }
    & & \\ \\

    \theutterance \stepcounter{utterance}  
    & & \multicolumn{4}{p{0.6\linewidth}}{
        \cellcolor[rgb]{0.9,0.9,0.9}{
            \makecell[{{p{\linewidth}}}]{
                \texttt{\tiny{[GM$\rangle$P2]}}
                \texttt{Συμμετέχετε σε ένα συνεργατικό παιχνίδι μαντεψιάς.} \\
\texttt{Ο στόχος είναι να βρείτε αν αυτό το πλέγμα και μια άλλη εικόνα που μόνο εγώ μπορώ να δω, είναι τα ίδια.} \\
\texttt{Παρακαλώ περιγράψτε πρώτα το πλέγμα σας. Στη συνέχεια, θα δώσω τη δική μου περιγραφή και μπορούμε να κάνουμε ο ένας στον άλλο ερωτήσεις σχετικά με τα πλέγματα για να καταλάβουμε αν είναι τα ίδια. Ξεκινήστε τώρα τη σύντομη περιγραφή του πλέγματός σας με "ΠΕΡΙΓΡΑΦΗ:" ακολουθούμενη από την περιγραφή. Μην προσθέσετε τίποτα άλλο. Αυτό είναι το πλέγμα σας:} \\
\texttt{X ▢ ▢ ▢ ▢} \\
\texttt{X ▢ ▢ ▢ ▢} \\
\texttt{X X X ▢ ▢} \\
\texttt{X ▢ ▢ ▢ ▢} \\
\texttt{X X X X X} \\
            }
        }
    }
    & \\ \\

    \theutterance \stepcounter{utterance}  
    & & & \multicolumn{4}{p{0.6\linewidth}}{
        \cellcolor[rgb]{0.9,0.9,0.9}{
            \makecell[{{p{\linewidth}}}]{
                \texttt{\tiny{[GM$\langle$P2]}}
                \texttt{ΠΕΡΙΓΡΑΦΗ: Ένα πλέγμα 5x5 όπου τα στοιχεία είναι είτε 'X' είτε '▢'. Η πρώτη στήλη έχει 'X' σε όλες τις θέσεις. Η τρίτη σειρά έχει 'X' στις τρεις πρώτες θέσεις. Η πέμπτη σειρά έχει 'X' σε όλες τις πέντε θέσεις. Οι υπόλοιπες θέσεις είναι '▢'.} \\
            }
        }
    }
     \\ \\

    \theutterance \stepcounter{utterance}  
    & & & \multicolumn{2}{p{0.3\linewidth}}{
        \cellcolor[rgb]{0.9,0.9,0.9}{
            \makecell[{{p{\linewidth}}}]{
                \texttt{\tiny{[GM$|$GM]}}
                \texttt{continue} \\
            }
        }
    }
    & & \\ \\

    \theutterance \stepcounter{utterance}  
    & & \multicolumn{4}{p{0.6\linewidth}}{
        \cellcolor[rgb]{0.9,0.9,0.9}{
            \makecell[{{p{\linewidth}}}]{
                \texttt{\tiny{[GM$\rangle$P2]}}
                \texttt{Αυτή είναι η ΠΕΡΙΓΡΑΦΗ μουΠΕΡΙΓΡΑΦΗ: Ένα πλέγμα 5x5 όπου οι πρώτες στήλες και οι πρώτες γραμμές περιέχουν κυρίως "X". Συγκεκριμένα, η πρώτη στήλη είναι γεμάτη με "X". Η πρώτη γραμμή είναι γεμάτη με "X". Η δεύτερη γραμμή έχει "X" στην πρώτη θέση και μετά τέσσερα κενά (▢). Η τρίτη γραμμή έχει τρία "X" στην αρχή και δύο κενά (▢). Οι τέσσερτες και πέμπτες γραμμές έχουν ένα "X" στην αρχή και τέσσερα κενά (▢).} \\
\texttt{Τώρα κάντε μια ερώτηση για να μάθετε νέες πτυχές του πλέγματός μου που μπορεί να διαφέρουν από το δικό σας πλέγμα. Ξεκινήστε με "ΕΡΩΤΗΣΗ:" και μην προσθέσετε τίποτα άλλο.} \\
            }
        }
    }
    & \\ \\

    \theutterance \stepcounter{utterance}  
    & & & \multicolumn{4}{p{0.6\linewidth}}{
        \cellcolor[rgb]{0.9,0.9,0.9}{
            \makecell[{{p{\linewidth}}}]{
                \texttt{\tiny{[GM$\langle$P2]}}
                \texttt{ΕΡΩΤΗΣΗ: Τι περιέχει η πέμπτη σειρά του δικού σας πλέγματος;} \\
            }
        }
    }
     \\ \\

    \theutterance \stepcounter{utterance}  
    & & & \multicolumn{2}{p{0.3\linewidth}}{
        \cellcolor[rgb]{0.9,0.9,0.9}{
            \makecell[{{p{\linewidth}}}]{
                \texttt{\tiny{[GM$|$GM]}}
                \texttt{continue} \\
            }
        }
    }
    & & \\ \\

    \theutterance \stepcounter{utterance}  
    & & \multicolumn{4}{p{0.6\linewidth}}{
        \cellcolor[rgb]{0.9,0.9,0.9}{
            \makecell[{{p{\linewidth}}}]{
                \texttt{\tiny{[P1$\langle$GM]}}
                \texttt{Αυτή είναι η ΠΕΡΙΓΡΑΦΗ μουΠΕΡΙΓΡΑΦΗ: Ένα πλέγμα 5x5 όπου τα στοιχεία είναι είτε 'X' είτε '▢'. Η πρώτη στήλη έχει 'X' σε όλες τις θέσεις. Η τρίτη σειρά έχει 'X' στις τρεις πρώτες θέσεις. Η πέμπτη σειρά έχει 'X' σε όλες τις πέντε θέσεις. Οι υπόλοιπες θέσεις είναι '▢'.} \\
\texttt{ΕΡΩΤΗΣΗ: Τι περιέχει η πέμπτη σειρά του δικού σας πλέγματος;Ξεκινήστε την απάντησή σας με "ΑΠΑΝΤΗΣΗ:" και μην προσθέσετε τίποτα άλλο.} \\
            }
        }
    }
    & \\ \\

    \theutterance \stepcounter{utterance}  
    & \multicolumn{4}{p{0.6\linewidth}}{
        \cellcolor[rgb]{0.9,0.9,0.9}{
            \makecell[{{p{\linewidth}}}]{
                \texttt{\tiny{[P1$\rangle$GM]}}
                \texttt{ΑΠΑΝΤΗΣΗ: Η πέμπτη σειρά του δικού μου πλέγματος περιέχει ένα "X" στην πρώτη θέση και τέσσερα κενά (▢) στις υπόλοιπες θέσεις (X ▢ ▢ ▢ ▢).} \\
            }
        }
    }
    & & \\ \\

    \theutterance \stepcounter{utterance}  
    & & & \multicolumn{2}{p{0.3\linewidth}}{
        \cellcolor[rgb]{0.9,0.9,0.9}{
            \makecell[{{p{\linewidth}}}]{
                \texttt{\tiny{[GM$|$GM]}}
                \texttt{continue} \\
            }
        }
    }
    & & \\ \\

    \theutterance \stepcounter{utterance}  
    & & \multicolumn{4}{p{0.6\linewidth}}{
        \cellcolor[rgb]{0.9,0.9,0.9}{
            \makecell[{{p{\linewidth}}}]{
                \texttt{\tiny{[P1$\langle$GM]}}
                \texttt{Τώρα κάντε μια ερώτηση για να μάθετε νέες πτυχές του πλέγματός μου που μπορεί να διαφέρουν από το δικό σας πλέγμα. Ξεκινήστε με "ΕΡΩΤΗΣΗ:" και μην προσθέσετε τίποτα άλλο.} \\
            }
        }
    }
    & \\ \\

    \theutterance \stepcounter{utterance}  
    & \multicolumn{4}{p{0.6\linewidth}}{
        \cellcolor[rgb]{0.9,0.9,0.9}{
            \makecell[{{p{\linewidth}}}]{
                \texttt{\tiny{[P1$\rangle$GM]}}
                \texttt{ΕΡΩΤΗΣΗ: Τι περιέχει η πρώτη γραμμή του δικού σας πλέγματος;} \\
            }
        }
    }
    & & \\ \\

    \theutterance \stepcounter{utterance}  
    & & & \multicolumn{2}{p{0.3\linewidth}}{
        \cellcolor[rgb]{0.9,0.9,0.9}{
            \makecell[{{p{\linewidth}}}]{
                \texttt{\tiny{[GM$|$GM]}}
                \texttt{continue} \\
            }
        }
    }
    & & \\ \\

    \theutterance \stepcounter{utterance}  
    & & \multicolumn{4}{p{0.6\linewidth}}{
        \cellcolor[rgb]{0.9,0.9,0.9}{
            \makecell[{{p{\linewidth}}}]{
                \texttt{\tiny{[GM$\rangle$P2]}}
                \texttt{ΑΠΑΝΤΗΣΗ: Η πέμπτη σειρά του δικού μου πλέγματος περιέχει ένα "X" στην πρώτη θέση και τέσσερα κενά (▢) στις υπόλοιπες θέσεις (X ▢ ▢ ▢ ▢).} \\
\texttt{ΕΡΩΤΗΣΗ: Τι περιέχει η πρώτη γραμμή του δικού σας πλέγματος;Ξεκινήστε την απάντησή σας με "ΑΠΑΝΤΗΣΗ:" και μην προσθέσετε τίποτα άλλο.} \\
            }
        }
    }
    & \\ \\

    \theutterance \stepcounter{utterance}  
    & & & \multicolumn{4}{p{0.6\linewidth}}{
        \cellcolor[rgb]{0.9,0.9,0.9}{
            \makecell[{{p{\linewidth}}}]{
                \texttt{\tiny{[GM$\langle$P2]}}
                \texttt{ΑΠΑΝΤΗΣΗ: Η πρώτη γραμμή του δικού μου πλέγματος περιέχει ένα "X" στην πρώτη θέση και τέσσερα κενά (▢) στις υπόλοιπες θέσεις (X ▢ ▢ ▢ ▢).} \\
            }
        }
    }
     \\ \\

    \theutterance \stepcounter{utterance}  
    & & & \multicolumn{2}{p{0.3\linewidth}}{
        \cellcolor[rgb]{0.9,0.9,0.9}{
            \makecell[{{p{\linewidth}}}]{
                \texttt{\tiny{[GM$|$GM]}}
                \texttt{continue} \\
            }
        }
    }
    & & \\ \\

    \theutterance \stepcounter{utterance}  
    & & \multicolumn{4}{p{0.6\linewidth}}{
        \cellcolor[rgb]{0.9,0.9,0.9}{
            \makecell[{{p{\linewidth}}}]{
                \texttt{\tiny{[GM$\rangle$P2]}}
                \texttt{Τώρα κάντε μια ερώτηση για να μάθετε νέες πτυχές του πλέγματός μου που μπορεί να διαφέρουν από το δικό σας πλέγμα. Ξεκινήστε με "ΕΡΩΤΗΣΗ:" και μην προσθέσετε τίποτα άλλο.} \\
            }
        }
    }
    & \\ \\

    \theutterance \stepcounter{utterance}  
    & & & \multicolumn{4}{p{0.6\linewidth}}{
        \cellcolor[rgb]{0.9,0.9,0.9}{
            \makecell[{{p{\linewidth}}}]{
                \texttt{\tiny{[GM$\langle$P2]}}
                \texttt{ΕΡΩΤΗΣΗ: Στο δικό σας πλέγμα, η πρώτη γραμμή περιέχει μόνο ένα "X" στην πρώτη θέση ή είναι γεμάτη με "X" (X X X X X);} \\
            }
        }
    }
     \\ \\

    \theutterance \stepcounter{utterance}  
    & & & \multicolumn{2}{p{0.3\linewidth}}{
        \cellcolor[rgb]{0.9,0.9,0.9}{
            \makecell[{{p{\linewidth}}}]{
                \texttt{\tiny{[GM$|$GM]}}
                \texttt{continue} \\
            }
        }
    }
    & & \\ \\

    \theutterance \stepcounter{utterance}  
    & & \multicolumn{4}{p{0.6\linewidth}}{
        \cellcolor[rgb]{0.9,0.9,0.9}{
            \makecell[{{p{\linewidth}}}]{
                \texttt{\tiny{[P1$\langle$GM]}}
                \texttt{ΑΠΑΝΤΗΣΗ: Η πρώτη γραμμή του δικού μου πλέγματος περιέχει ένα "X" στην πρώτη θέση και τέσσερα κενά (▢) στις υπόλοιπες θέσεις (X ▢ ▢ ▢ ▢).} \\
\texttt{ΕΡΩΤΗΣΗ: Στο δικό σας πλέγμα, η πρώτη γραμμή περιέχει μόνο ένα "X" στην πρώτη θέση ή είναι γεμάτη με "X" (X X X X X);Ξεκινήστε την απάντησή σας με "ΑΠΑΝΤΗΣΗ:" και μην προσθέσετε τίποτα άλλο.} \\
            }
        }
    }
    & \\ \\

    \theutterance \stepcounter{utterance}  
    & \multicolumn{4}{p{0.6\linewidth}}{
        \cellcolor[rgb]{0.9,0.9,0.9}{
            \makecell[{{p{\linewidth}}}]{
                \texttt{\tiny{[P1$\rangle$GM]}}
                \texttt{ΑΠΑΝΤΗΣΗ: Η πρώτη γραμμή του δικού μου πλέγματος είναι γεμάτη με "X" (X X X X X).} \\
            }
        }
    }
    & & \\ \\

    \theutterance \stepcounter{utterance}  
    & & & \multicolumn{2}{p{0.3\linewidth}}{
        \cellcolor[rgb]{0.9,0.9,0.9}{
            \makecell[{{p{\linewidth}}}]{
                \texttt{\tiny{[GM$|$GM]}}
                \texttt{continue} \\
            }
        }
    }
    & & \\ \\

    \theutterance \stepcounter{utterance}  
    & & \multicolumn{4}{p{0.6\linewidth}}{
        \cellcolor[rgb]{0.9,0.9,0.9}{
            \makecell[{{p{\linewidth}}}]{
                \texttt{\tiny{[P1$\langle$GM]}}
                \texttt{Τώρα κάντε μια ερώτηση για να μάθετε νέες πτυχές του πλέγματός μου που μπορεί να διαφέρουν από το δικό σας πλέγμα. Ξεκινήστε με "ΕΡΩΤΗΣΗ:" και μην προσθέσετε τίποτα άλλο.} \\
            }
        }
    }
    & \\ \\

    \theutterance \stepcounter{utterance}  
    & \multicolumn{4}{p{0.6\linewidth}}{
        \cellcolor[rgb]{0.9,0.9,0.9}{
            \makecell[{{p{\linewidth}}}]{
                \texttt{\tiny{[P1$\rangle$GM]}}
                \texttt{ΕΡΩΤΗΣΗ: Τι περιέχει η δεύτερη γραμμή του δικού σας πλέγματος;} \\
            }
        }
    }
    & & \\ \\

    \theutterance \stepcounter{utterance}  
    & & & \multicolumn{2}{p{0.3\linewidth}}{
        \cellcolor[rgb]{0.9,0.9,0.9}{
            \makecell[{{p{\linewidth}}}]{
                \texttt{\tiny{[GM$|$GM]}}
                \texttt{continue} \\
            }
        }
    }
    & & \\ \\

    \theutterance \stepcounter{utterance}  
    & & \multicolumn{4}{p{0.6\linewidth}}{
        \cellcolor[rgb]{0.9,0.9,0.9}{
            \makecell[{{p{\linewidth}}}]{
                \texttt{\tiny{[GM$\rangle$P2]}}
                \texttt{ΑΠΑΝΤΗΣΗ: Η πρώτη γραμμή του δικού μου πλέγματος είναι γεμάτη με "X" (X X X X X).} \\
\texttt{ΕΡΩΤΗΣΗ: Τι περιέχει η δεύτερη γραμμή του δικού σας πλέγματος;Ξεκινήστε την απάντησή σας με "ΑΠΑΝΤΗΣΗ:" και μην προσθέσετε τίποτα άλλο.} \\
            }
        }
    }
    & \\ \\

    \theutterance \stepcounter{utterance}  
    & & & \multicolumn{4}{p{0.6\linewidth}}{
        \cellcolor[rgb]{0.9,0.9,0.9}{
            \makecell[{{p{\linewidth}}}]{
                \texttt{\tiny{[GM$\langle$P2]}}
                \texttt{ΑΠΑΝΤΗΣΗ: Η δεύτερη γραμμή του δικού μου πλέγματος περιέχει ένα "X" στην πρώτη θέση και τέσσερα κενά (▢) στις υπόλοιπες θέσεις (X ▢ ▢ ▢ ▢).} \\
            }
        }
    }
     \\ \\

    \theutterance \stepcounter{utterance}  
    & & & \multicolumn{2}{p{0.3\linewidth}}{
        \cellcolor[rgb]{0.9,0.9,0.9}{
            \makecell[{{p{\linewidth}}}]{
                \texttt{\tiny{[GM$|$GM]}}
                \texttt{continue} \\
            }
        }
    }
    & & \\ \\

    \theutterance \stepcounter{utterance}  
    & & \multicolumn{4}{p{0.6\linewidth}}{
        \cellcolor[rgb]{0.9,0.9,0.9}{
            \makecell[{{p{\linewidth}}}]{
                \texttt{\tiny{[GM$\rangle$P2]}}
                \texttt{Τώρα κάντε μια ερώτηση για να μάθετε νέες πτυχές του πλέγματός μου που μπορεί να διαφέρουν από το δικό σας πλέγμα. Ξεκινήστε με "ΕΡΩΤΗΣΗ:" και μην προσθέσετε τίποτα άλλο.} \\
            }
        }
    }
    & \\ \\

    \theutterance \stepcounter{utterance}  
    & & & \multicolumn{4}{p{0.6\linewidth}}{
        \cellcolor[rgb]{0.9,0.9,0.9}{
            \makecell[{{p{\linewidth}}}]{
                \texttt{\tiny{[GM$\langle$P2]}}
                \texttt{ΕΡΩΤΗΣΗ: Στο δικό σας πλέγμα, η πέμπτη σειρά περιέχει πέντε "X" (X X X X X) ή μόνο ένα "X" στην πρώτη θέση (X ▢ ▢ ▢ ▢);} \\
            }
        }
    }
     \\ \\

    \theutterance \stepcounter{utterance}  
    & & & \multicolumn{2}{p{0.3\linewidth}}{
        \cellcolor[rgb]{0.9,0.9,0.9}{
            \makecell[{{p{\linewidth}}}]{
                \texttt{\tiny{[GM$|$GM]}}
                \texttt{continue} \\
            }
        }
    }
    & & \\ \\

    \theutterance \stepcounter{utterance}  
    & & \multicolumn{4}{p{0.6\linewidth}}{
        \cellcolor[rgb]{0.9,0.9,0.9}{
            \makecell[{{p{\linewidth}}}]{
                \texttt{\tiny{[P1$\langle$GM]}}
                \texttt{ΑΠΑΝΤΗΣΗ: Η δεύτερη γραμμή του δικού μου πλέγματος περιέχει ένα "X" στην πρώτη θέση και τέσσερα κενά (▢) στις υπόλοιπες θέσεις (X ▢ ▢ ▢ ▢).} \\
\texttt{ΕΡΩΤΗΣΗ: Στο δικό σας πλέγμα, η πέμπτη σειρά περιέχει πέντε "X" (X X X X X) ή μόνο ένα "X" στην πρώτη θέση (X ▢ ▢ ▢ ▢);Ξεκινήστε την απάντησή σας με "ΑΠΑΝΤΗΣΗ:" και μην προσθέσετε τίποτα άλλο.} \\
            }
        }
    }
    & \\ \\

    \theutterance \stepcounter{utterance}  
    & \multicolumn{4}{p{0.6\linewidth}}{
        \cellcolor[rgb]{0.9,0.9,0.9}{
            \makecell[{{p{\linewidth}}}]{
                \texttt{\tiny{[P1$\rangle$GM]}}
                \texttt{ΑΠΑΝΤΗΣΗ: Η πέμπτη σειρά του δικού μου πλέγματος περιέχει μόνο ένα "X" στην πρώτη θέση (X ▢ ▢ ▢ ▢).} \\
            }
        }
    }
    & & \\ \\

    \theutterance \stepcounter{utterance}  
    & & & \multicolumn{2}{p{0.3\linewidth}}{
        \cellcolor[rgb]{0.9,0.9,0.9}{
            \makecell[{{p{\linewidth}}}]{
                \texttt{\tiny{[GM$|$GM]}}
                \texttt{continue} \\
            }
        }
    }
    & & \\ \\

    \theutterance \stepcounter{utterance}  
    & & \multicolumn{4}{p{0.6\linewidth}}{
        \cellcolor[rgb]{0.9,0.9,0.9}{
            \makecell[{{p{\linewidth}}}]{
                \texttt{\tiny{[P1$\langle$GM]}}
                \texttt{Τώρα κάντε μια ερώτηση για να μάθετε νέες πτυχές του πλέγματός μου που μπορεί να διαφέρουν από το δικό σας πλέγμα. Ξεκινήστε με "ΕΡΩΤΗΣΗ:" και μην προσθέσετε τίποτα άλλο.} \\
            }
        }
    }
    & \\ \\

    \theutterance \stepcounter{utterance}  
    & \multicolumn{4}{p{0.6\linewidth}}{
        \cellcolor[rgb]{0.9,0.9,0.9}{
            \makecell[{{p{\linewidth}}}]{
                \texttt{\tiny{[P1$\rangle$GM]}}
                \texttt{ΕΡΩΤΗΣΗ: Τι περιέχει η δεύτερη γραμμή του δικού σας πλέγματος;} \\
            }
        }
    }
    & & \\ \\

    \theutterance \stepcounter{utterance}  
    & & & \multicolumn{2}{p{0.3\linewidth}}{
        \cellcolor[rgb]{0.9,0.9,0.9}{
            \makecell[{{p{\linewidth}}}]{
                \texttt{\tiny{[GM$|$GM]}}
                \texttt{continue} \\
            }
        }
    }
    & & \\ \\

    \theutterance \stepcounter{utterance}  
    & & \multicolumn{4}{p{0.6\linewidth}}{
        \cellcolor[rgb]{0.9,0.9,0.9}{
            \makecell[{{p{\linewidth}}}]{
                \texttt{\tiny{[GM$\rangle$P2]}}
                \texttt{ΑΠΑΝΤΗΣΗ: Η πέμπτη σειρά του δικού μου πλέγματος περιέχει μόνο ένα "X" στην πρώτη θέση (X ▢ ▢ ▢ ▢).} \\
\texttt{ΕΡΩΤΗΣΗ: Τι περιέχει η δεύτερη γραμμή του δικού σας πλέγματος;Ξεκινήστε την απάντησή σας με "ΑΠΑΝΤΗΣΗ:" και μην προσθέσετε τίποτα άλλο.} \\
            }
        }
    }
    & \\ \\

    \theutterance \stepcounter{utterance}  
    & & & \multicolumn{4}{p{0.6\linewidth}}{
        \cellcolor[rgb]{0.9,0.9,0.9}{
            \makecell[{{p{\linewidth}}}]{
                \texttt{\tiny{[GM$\langle$P2]}}
                \texttt{ΑΠΑΝΤΗΣΗ: Η δεύτερη γραμμή του δικού μου πλέγματος περιέχει ένα "X" στην πρώτη θέση και τέσσερα κενά (▢) στις υπόλοιπες θέσεις (X ▢ ▢ ▢ ▢).} \\
            }
        }
    }
     \\ \\

    \theutterance \stepcounter{utterance}  
    & & & \multicolumn{2}{p{0.3\linewidth}}{
        \cellcolor[rgb]{0.9,0.9,0.9}{
            \makecell[{{p{\linewidth}}}]{
                \texttt{\tiny{[GM$|$GM]}}
                \texttt{continue} \\
            }
        }
    }
    & & \\ \\

    \theutterance \stepcounter{utterance}  
    & & \multicolumn{4}{p{0.6\linewidth}}{
        \cellcolor[rgb]{0.9,0.9,0.9}{
            \makecell[{{p{\linewidth}}}]{
                \texttt{\tiny{[GM$\rangle$P2]}}
                \texttt{Τώρα πάρε μια απόφαση. Τι πιστεύεις: είναι το δικό σου πλέγμα και το άλλο πλέγμα που περιγράφεται το ίδιο; Γράψε "ΑΠΟΦΑΣΗ: ίδιο πλέγμα" αν πιστεύεις ότι είναι το ίδιο πλέγμα ή "ΑΠΟΦΑΣΗ: διαφορετικά πλέγματα" αν πιστεύεις ότι είναι διαφορετικά πλέγματα. Μην προσθέσεις τίποτα άλλο.} \\
            }
        }
    }
    & \\ \\

    \theutterance \stepcounter{utterance}  
    & & & \multicolumn{4}{p{0.6\linewidth}}{
        \cellcolor[rgb]{0.9,0.9,0.9}{
            \makecell[{{p{\linewidth}}}]{
                \texttt{\tiny{[GM$\langle$P2]}}
                \texttt{ΑΠΟΦΑΣΗ: διαφορετικά πλέγματα} \\
            }
        }
    }
     \\ \\

    \theutterance \stepcounter{utterance}  
    & & & \multicolumn{2}{p{0.3\linewidth}}{
        \cellcolor[rgb]{0.9,0.9,0.9}{
            \makecell[{{p{\linewidth}}}]{
                \texttt{\tiny{[GM$|$GM]}}
                \texttt{continue} \\
            }
        }
    }
    & & \\ \\

    \theutterance \stepcounter{utterance}  
    & & & \multicolumn{2}{p{0.3\linewidth}}{
        \cellcolor[rgb]{0.9,0.9,0.9}{
            \makecell[{{p{\linewidth}}}]{
                \texttt{\tiny{[GM$|$GM]}}
                \texttt{success} \\
            }
        }
    }
    & & \\ \\

    \theutterance \stepcounter{utterance}  
    & & \multicolumn{4}{p{0.6\linewidth}}{
        \cellcolor[rgb]{0.9,0.9,0.9}{
            \makecell[{{p{\linewidth}}}]{
                \texttt{\tiny{[P1$\langle$GM]}}
                \texttt{ΑΠΑΝΤΗΣΗ: Η δεύτερη γραμμή του δικού μου πλέγματος περιέχει ένα "X" στην πρώτη θέση και τέσσερα κενά (▢) στις υπόλοιπες θέσεις (X ▢ ▢ ▢ ▢).} \\
\texttt{Τώρα πάρε μια απόφαση. Τι πιστεύεις: είναι το δικό σου πλέγμα και το άλλο πλέγμα που περιγράφεται το ίδιο; Γράψε "ΑΠΟΦΑΣΗ: ίδιο πλέγμα" αν πιστεύεις ότι είναι το ίδιο πλέγμα ή "ΑΠΟΦΑΣΗ: διαφορετικά πλέγματα" αν πιστεύεις ότι είναι διαφορετικά πλέγματα. Μην προσθέσεις τίποτα άλλο.} \\
            }
        }
    }
    & \\ \\

    \theutterance \stepcounter{utterance}  
    & \multicolumn{4}{p{0.6\linewidth}}{
        \cellcolor[rgb]{0.9,0.9,0.9}{
            \makecell[{{p{\linewidth}}}]{
                \texttt{\tiny{[P1$\rangle$GM]}}
                \texttt{ΑΠΟΦΑΣΗ: διαφορετικά πλέγματα} \\
            }
        }
    }
    & & \\ \\

    \theutterance \stepcounter{utterance}  
    & & & \multicolumn{2}{p{0.3\linewidth}}{
        \cellcolor[rgb]{0.9,0.9,0.9}{
            \makecell[{{p{\linewidth}}}]{
                \texttt{\tiny{[GM$|$GM]}}
                \texttt{continue} \\
            }
        }
    }
    & & \\ \\

    \theutterance \stepcounter{utterance}  
    & & & \multicolumn{2}{p{0.3\linewidth}}{
        \cellcolor[rgb]{0.9,0.9,0.9}{
            \makecell[{{p{\linewidth}}}]{
                \texttt{\tiny{[GM$|$GM]}}
                \texttt{success} \\
            }
        }
    }
    & & \\ \\

\end{supertabular}
}

\end{document}
